Krátké shrnutí problému
Detektory „AI‑textu“ mohou pomoci, ale nejsou spolehlivé ani právně/eticky bezproblémové; níže stručně proč a praktické alternativy.
Proč detektorům nevěřit \begin{itemize}
  \item Nejsou definitivním důkazem: falešné pozitivy/negativy; neřeší autorství ani kontext.
  \item Jako hlavní nástroj mohou vést k nespravedlnosti a poškození důvěry.
\end{itemize}
Praktické alternativy a doporučené postupy
1) Požadavek na procesní artefakty \begin{itemize}
  \item Požadujte meziverze, commit logy, exporty nebo screenshoty; hodnocení se zaměří na proces \cite{kb_ai_101_tipu_pdf_04e4a115_page_8_1}.
  \item V LMS mějte povinné pole pro průběžné verze a šablony.
\end{itemize}
2) Úlohy odolné vůči triviálnímu použití AI \begin{itemize}
  \item Dejte přednost kontextovým úlohám, místním zdrojům a iteracím; používejte časově omezená in‑class writings, ústní zkoušky či portfolia.
\end{itemize}
3) Transparentní deklarace a reflexe použití AI \begin{itemize}
  \item Pokud je AI povolena, vyžadujte deklaraci a krátkou sebereflexi o roli nástroje a ověření výstupů.
\end{itemize}
4) Ověřovací rozhovor \begin{itemize}
  \item Při podezření doplňte krátký pohovor, kde student vysvětlí postup a rozhodnutí.
\end{itemize}
5) Školní/enterprise řešení \begin{itemize}
  \item Licencované balíky nabízejí lepší kontrolu dat, auditní logy a integraci do výuky \cite{kb_ai_101_tipu_pdf_04e4a115_page_36_1,kb_ai_101_tipu_pdf_04e4a115_page_15_2}.
\end{itemize}
Postup při podezření na nevhodné použití AI \begin{enumerate}
  \item Zkontrolujte povinné procesní artefakty; chybějící vyžádejte \cite{kb_ai_101_tipu_pdf_04e4a115_page_8_1}.
  \item V pohovoru se zaměřte na vysvětlení postupu, zdrojů a rozhodnutí.
  \item Dokumentujte zjištění podle interních pravidel.
  \item Pravidelně vyhodnocujte a upravujte opatření.
\end{enumerate}
Závěrečná poznámka
Místo spoléhání se na automaty doporučujeme kombinaci dobře navržených úloh, povinných artefaktů, transparentních pravidel a krátkých ověřovacích pohovorů; tyto přístupy lépe chrání akademickou integritu a respektují právní a etické aspekty práce se studentskými daty \cite{kb_ai_101_tipu_pdf_04e4a115_page_8_1,kb_ai_101_tipu_pdf_04e4a115_page_36_1,kb_ai_101_tipu_pdf_04e4a115_page_15_2}.